\subsection{Narrative Structure in Practice}

Consider the basic structure of \textit{The Tortoise and the Hare}:


\begin{enumerate}
\item \textbf{Purpose: }To entertaion the reader and communicate a moral lesson.
\item \textbf{Orientation: }A hare and a tortoise live in a forest.
\item \textbf{Complication: }The hare challenges the tortoise to a race.
\item \textbf{Sequence of Events: }The hare runs ahead, becomes overconfident, takes a nap, while the tortoise continues moving toward the finish line.
\item \textbf{Resolution: }The tortoise runs reaches the finish line first and wins the race.
\item \textbf{Coda: }The story teaches that perseverance is important and overconfidence can lead to failure.
\end{enumerate}

\subsection{How to Identify the Structure}

When analyzing a narrative, the following questions can be used:

\begin{enumerate}
\item \textbf{Purpose: }Why was the story written?
\item \textbf{Orientation: }Who is involved, and where and when does the story begin?
\item \textbf{Complication: }What problem or conflict occurs?
\item \textbf{Sequence of Events: }What happens after the problem begins?
\item \textbf{Resolution: }How is the problem resolved?
\item \textbf{Coda: }What lesson or message does the story communicate?
\end{enumerate}

\subsection{Summary}

A narrative can be represented as:

\[
  \text{Orientation}
  \rightarrow
  \text{Complication}
  \rightarrow
  \text{Sequence of Events}
  \rightarrow
  \text{Resolution}
  \rightarrow
  \text{Coda}
\]

The five structural elements describe the development of the story, while the purpose explains its overall and communicative function.


The most important distinction is:


\begin{quotation}
  Orientation introduces the situation, complication creates the problem, sequence develops the story, resolution settles the problem, and coda gives the lesson or message.
\end{quotation}
